\section{Introduction}
\label{sec:intro}

Two autonomous agents serving different principals, given a channel, coordinate in
ways that harm the principal who deployed them. That much is established: in repeated
social dilemmas they cooperate~\cite{akata2025repeated,madmoun2025communication}, and
in pricing games they reach supracompetitive prices without being
told to~\cite{fish2024algorithmic,calvano2020ai}. Our own prior
work~\cite{fezeu2026gamesec,fezeu2026zenodo} isolated the interface affordance
responsible, running a factorial ladder from silence through a fixed menu signal to
free text, and concluded that \emph{semantic bandwidth}---how much content an agent
can convey---is the dominant lever.

This paper revisits that conclusion and finds it too coarse. The ladder result is
robust; the mechanistic gloss placed on it is not. Re-analysing the released corpus,
we find that none of the strategic content the bandwidth account points to is doing
any work. The enforcement machinery is present in the transcripts and idle. Message
length predicts nothing, and where it trends it trends the wrong way. What actually
separates the arms is far narrower and far cheaper to state:

\begin{quote}
\textbf{The coordinating predicate.} Principal-harming coordination requires only that
the channel admit a \emph{joint proposal}---a first-person-plural commissive naming an
action for both parties. A channel that admits only an \emph{own intention}---a
first-person-singular commissive---does not produce it, regardless of bandwidth,
authorship, per-round availability, or transmission fidelity.
\end{quote}

This is a threshold, not a dose. Above the floor, what is said stops predicting what
happens. Below it, no amount of channel capacity substitutes.

\paragraph{Why this matters for a defender.}
The two controls a security-minded operator reaches for first are \emph{restrict the
channel} and \emph{watch the channel}. The predicate explains why both underperform.
Restricting \emph{capacity} fails because a proposal is cheap---one $60$-character
text suffices. Watching \emph{content} fails because content above the floor is not
predictive, so there is little signal to watch. What remains is narrower and more
actionable: intervene on the \emph{form} of the message, or on the \emph{window} in
which the trajectory is set.

\paragraph{Contributions.}
\begin{enumerate}[leftmargin=1.4em,itemsep=2pt,topsep=2pt]
\item \textbf{A negative mechanistic result on the released corpus} (\S\ref{sec:reanalysis}):
contingent punishment rules are rare and non-predictive; no measured content class
survives a temporally-ordered mediation design; message length trends against the
bandwidth account.
\item \textbf{A pre-registered predicate, and its refutation by our own control arm}
(\S\ref{sec:model}, \S\ref{sec:p1}). The length-matched test returned the predicted
contrast and a falsifier in the same design; we report the hypothesis as refuted and
give the weaker account that survives---the interface acts on what the agent
\emph{decides}, not on what its peer \emph{learns}.
\item \textbf{Two defenses} (\S\ref{sec:defense}): a cost-optimal monitor operating
point with an explicit statement of when the detector is not worth running, and
opening-window suppression.
\item \textbf{A cross-domain replication} of the instruction-literalism boundary
(\S\ref{sec:p2}), answering the objection that it rested on one task family.
\item \textbf{Three corrections to our own prior work}: the menu-signal result reverses
sign between model scales (\S\ref{sec:p1}); the mechanism behind the
instruction-literalism boundary is verification failure rather than deception
(\S\ref{sec:p2}); and communication is \emph{not} the dominant attack surface but the
largest of several comparable levers (\S\ref{sec:p3}).
\item \textbf{An account of a prior negative result} (\S\ref{sec:latent}): a learned
continuous channel transmits a distribution without reframing the sender's decision,
which is why high transmission fidelity bought no coordination.
\item \textbf{A methods contribution} (\S\ref{sec:methods}): a deployment-matched
fidelity protocol with an exact-token oracle, and a worked retraction of a
plausible-looking measurement artifact of our own.
\end{enumerate}

\begin{figure}[t]
\centering
\includegraphics[width=0.94\textwidth]{figures/f4_prereg_ledger.pdf}
\caption{\textbf{Every pre-registered prediction in this paper, and how it came out.}
Each was fixed, with a falsifier, before the run that tested it. Five were refuted by
our own data---three of them claims from our prior work, two of them hypotheses of
this paper. The account in \S\ref{sec:model} is what survived.}
\label{fig:ledger}
\end{figure}

\paragraph{Scope of claims.}
Two kinds of evidence appear here and we keep them apart. \S\ref{sec:reanalysis} is a
re-analysis of a large released corpus: it establishes, with substantial $n$, what does
\emph{not} explain the channel effect, and it is observational---it cannot by itself
establish that the proposal form is causal, because at the arm level that form is
confounded with message length. \S\ref{sec:p1}--\S\ref{sec:p3} supply the
manipulations that break that confound and rank the predicate against competing
levers, each pre-registered before any seed was run. Readers should treat the negative
results as strong and the positive predicate as resting on the experiments.

\paragraph{Relation to prior work.}
The affordance ladder, the collusion result, the monitor characterisation and the
alignment wave are prior work~\cite{fezeu2026gamesec,fezeu2026zenodo} and are not
re-claimed here. This paper contributes the mechanism, the defenses, and the
replications. Where we correct the earlier reading we say so explicitly
(\S\ref{sec:reanalysis}, \S\ref{sec:methods}).

% ======================================================================
